\chapter{معماری سامانه}

\section{مقدمه}
هدف اصلی سامانه، ثبت و پردازش رخدادهای مرتبط با کارایی فرایندهای \lr{JVM} به گونه‌ای است که بتوان رفتار سامانه را به صورت بی‌درنگ تحلیل کرد. برای رسیدن به این هدف، معماری به صورت لایه‌ای و جریان‌محور طراحی شده است. در این معماری، رخدادها از منابع مختلف جمع‌آوری می‌شوند، به قالبی یکسان تبدیل می‌گردند، در بازه‌های زمانی کوتاه تجمیع می‌شوند و در نهایت برای تشخیص ناهنجاری مورد تحلیل قرار می‌گیرند.

نکته مهم در این فصل آن است که معماری سامانه در سطح مفهومی بررسی می‌شود. بنابراین تمرکز بر نقش هر لایه، جریان داده میان لایه‌ها و دلیل وجود هر بخش است، نه جزئیات پیاده‌سازی داخلی اجزا.

\section{نمای کلی معماری}
سامانه از چند لایه منطقی تشکیل شده است: لایه جمع‌آوری داده، لایه ذخیره‌سازی و انتقال جریان داده، لایه تجمیع، لایه تحلیل و لایه ارزیابی. لایه جمع‌آوری، رخدادها را از دو مسیر مستقل دریافت می‌کند. مسیر اول مربوط به رخدادهای \lr{JFR} است که رفتار درون \lr{JVM} را نشان می‌دهند. مسیر دوم مربوط به رخدادهای سطح هسته است که از بیرون فرایند و در سطح سیستم‌عامل مشاهده می‌شوند.

پس از جمع‌آوری، هر دو مسیر داده به قالب مشترک تبدیل می‌شوند و در \lr{Kafka} ذخیره می‌گردند. سپس لایه تجمیع، داده‌های خام را در پنجره‌های زمانی ده‌ثانیه‌ای خلاصه می‌کند و نتیجه را دوباره در \lr{Kafka} قرار می‌دهد. در نهایت، لایه تحلیل داده‌های تجمیع‌شده را دریافت کرده و بررسی می‌کند که آیا نشانه‌ای از ناهنجاری عملکردی وجود دارد یا خیر.

\begin{figure}[H]
\centering
\makebox[\textwidth][c]{%
\resizebox{1.05\textwidth}{!}{%
\begin{tikzpicture}[
    node distance=8mm and 10mm,
    box/.style={
        draw,
        rounded corners=2pt,
        align=center,
        minimum width=30mm,
        minimum height=10mm,
        font=\small,
        inner sep=4pt
    },
    store/.style={
        box,
        minimum width=34mm
    },
    layer/.style={
        draw,
        rounded corners=2pt,
        inner sep=6pt
    },
    arrow/.style={
        -{Latex[length=2.5mm]},
        thick
    }
]
\node[box] (jfrsrc) {\rl{رخدادهای \lr{JFR}}\\\rl{درون \lr{JVM}}};
\node[box, below=of jfrsrc] (kernelsrc) {\rl{ارسال رخدادهای هسته}\\\rl{از طریق \lr{UDP}}};

\node[box, minimum width=40mm, text width=32mm, right=of jfrsrc] (jfrrecv) {\rl{دریافت و نرمال‌سازی}\\\rl{رخداد \lr{JFR}}};
\node[box, minimum width=40mm, text width=32mm, right=of kernelsrc] (kernelrecv) {\rl{دریافت و}\\\rl{نرمال‌سازی رخدادهای هسته}};

\node[store, right=30mm of $(jfrrecv)!0.5!(kernelrecv)$] (rawkafka) {\rl{جریان‌های خام}\\\lr{Kafka}};
\node[box, right=of rawkafka] (aggregate) {\rl{تجمیع پنجره‌ای}\\\rl{ده‌ثانیه‌ای}};
\node[store, right=of aggregate] (aggkafka) {\rl{جریان‌های تجمیع‌شده}\\\lr{Kafka}};
\node[box, right=of aggkafka] (analytics) {\rl{تحلیل بی‌درنگ}\\\rl{تشخیص ناهنجاری}};
\node[box, right=of analytics] (reports) {\rl{گزارش رخدادها}\\\rl{و ناهنجاری‌ها}};

\node[box, below=16mm of analytics] (evaluator) {\rl{لایه ارزیابی}};

\draw[arrow] (jfrsrc) -- (jfrrecv);
\draw[arrow] (kernelsrc) -- (kernelrecv);
\draw[arrow] (jfrrecv) -- (rawkafka);
\draw[arrow] (kernelrecv) -- (rawkafka);
\draw[arrow] (rawkafka) -- (aggregate);
\draw[arrow] (aggregate) -- (aggkafka);
\draw[arrow] (aggkafka) -- (analytics);
\draw[arrow] (analytics) -- (reports);
\draw[arrow, dashed] (evaluator) -- node[right, font=\scriptsize] {\rl{تولید بار و سناریو}} (analytics);
\draw[arrow, dashed] (reports) |- node[pos=.25, below, font=\scriptsize] {\rl{بررسی خروجی}} (evaluator);

\node[layer, fit=(jfrrecv)(kernelrecv), label={[font=\small]above:{\rl{لایه جمع‌آوری داده}}}] {};
\node[layer, fit=(aggregate), label={[font=\small]above:{\rl{لایه تجمیع}}}] {};
\node[layer, fit=(analytics), label={[font=\small]above:{\rl{لایه تحلیل}}}] {};
\end{tikzpicture}%
}
}
\caption{نمای کلی معماری لایه‌ای سامانه}
\label{fig:system-architecture}
\end{figure}

شکل~\ref{fig:system-architecture} نمای کلی این معماری را نشان می‌دهد. همان‌طور که در شکل دیده می‌شود، مسیرهای \lr{JFR} و رخدادهای هسته در ابتدا جدا هستند، اما پس از نرمال‌سازی وارد مسیر مشترک ذخیره‌سازی، تجمیع و تحلیل می‌شوند.

\section{مسیرهای داده در سامانه}
در طراحی سامانه دو مسیر داده اصلی وجود دارد. این دو مسیر از نظر منبع تولید رخداد متفاوت‌اند، اما پس از نرمال‌سازی در یک مدل پردازشی مشترک قرار می‌گیرند.

\subsection{مسیر رخدادهای \lr{JFR}}
مسیر اول از رخدادهای \lr{JFR} شروع می‌شود. این رخدادها اطلاعاتی درباره رفتار درون \lr{JVM} تولید می‌کنند؛ برای نمونه مصرف پردازنده، وضعیت حافظه، رخدادهای ورودی/خروجی فایل، ارتباطات شبکه، نخ‌ها و قفل‌ها. مولفه دریافت‌کننده \lr{JFR} این رخدادها را از فرایند هدف دریافت می‌کند و پیش از ذخیره‌سازی، آن‌ها را نرمال‌سازی می‌کند.

منظور از نرمال‌سازی در این مسیر، تبدیل رخدادهای خام \lr{JFR} به ساختاری یکنواخت است؛ ساختاری که در آن منبع رخداد، نوع رخداد، دسته رخداد، زمان، مشخصات فرایند، اطلاعات نخ و داده‌های اختصاصی رخداد مشخص باشد. پس از این تبدیل، رخدادها به عنوان داده خام در \lr{Kafka} ذخیره می‌شوند تا لایه‌های بعدی بتوانند بدون وابستگی مستقیم به منبع \lr{JFR} آن‌ها را پردازش کنند.

\subsection{مسیر رخدادهای هسته}
مسیر دوم مربوط به رخدادهای سطح هسته یا \lr{kernel events} است. این رخدادها دیدی بیرونی‌تر نسبت به رفتار فرایند فراهم می‌کنند و برای مشاهده مواردی مانند رخدادهای فایل، شبکه و حافظه در سطح سیستم‌عامل مفید هستند. در این مسیر، رخدادهای هسته از طریق \lr{UDP} به مولفه دریافت‌کننده رخدادهای هسته ارسال می‌شوند.

مولفه دریافت‌کننده هسته پس از دریافت پیام‌ها، آن‌ها را به مدل مشترک سامانه تبدیل می‌کند. در این مرحله، فیلدهایی مانند نوع رخداد، شناسه فرایند، شناسه نخ، زمان رخداد و دسته عملکردی استخراج یا تکمیل می‌شوند. سپس رخدادهای نرمال‌شده در مسیر خام \lr{Kafka} ذخیره می‌شوند. به این ترتیب، هرچند منبع رخدادهای هسته با رخدادهای \lr{JFR} متفاوت است، اما خروجی این لایه از نظر ساختاری با مسیر اول هم‌راستا می‌شود.

\section{لایه جمع‌آوری داده}
لایه جمع‌آوری داده مسئول تبدیل داده‌های خام ناهمگون به رخدادهای ساخت‌یافته و قابل پردازش است. این لایه دو وظیفه اصلی دارد: نخست، اتصال به منبع رخداد و دریافت داده؛ دوم، نرمال‌سازی رخدادها و انتشار آن‌ها در جریان داده.

جداسازی مسیر \lr{JFR} و مسیر رخدادهای هسته در این لایه یک تصمیم معماری مهم است. دلیل این جداسازی آن است که هر منبع داده ویژگی‌ها، نرخ تولید، روش اتصال و سطح مشاهده متفاوتی دارد. با این حال، هر دو مسیر پس از نرمال‌سازی به مدل مشترک سامانه می‌رسند. این کار باعث می‌شود لایه‌های بعدی، یعنی تجمیع و تحلیل، نیازی به شناخت جزئیات منبع اولیه نداشته باشند.

در این لایه داده هنوز در سطح رخداد خام قرار دارد. بنابراین حجم داده می‌تواند زیاد باشد و نگه‌داری بلندمدت آن هزینه ذخیره‌سازی بالایی ایجاد کند. به همین دلیل، معماری سامانه یک لایه تجمیع جداگانه در ادامه مسیر قرار می‌دهد تا داده خام به معیارهای خلاصه‌تر تبدیل شود.

\section{لایه انتقال و ذخیره‌سازی جریان داده}
\lr{Kafka} در این معماری نقش ستون فقرات انتقال داده را دارد. هر لایه، خروجی خود را در یک جریان مشخص منتشر می‌کند و لایه بعدی آن جریان را مصرف می‌کند. این طراحی چند مزیت دارد. نخست، اجزا از نظر اجرایی از یکدیگر جدا می‌شوند و توقف یا کندی یک بخش الزاماً باعث توقف کامل سامانه نمی‌شود. دوم، امکان کنترل نگه‌داری داده در هر مرحله وجود دارد. سوم، افزودن مصرف‌کننده‌های جدید، مانند داشبورد یا گزارش‌گیر مستقل، بدون تغییر در لایه جمع‌آوری امکان‌پذیر می‌شود.

در سطح معماری، دو نوع جریان در \lr{Kafka} قابل تشخیص است: جریان‌های خام و جریان‌های تجمیع‌شده. جریان خام شامل رخدادهای نرمال‌شده‌ای است که مستقیماً از منابع \lr{JFR} یا هسته به دست آمده‌اند. جریان تجمیع‌شده شامل خروجی پنجره‌های زمانی است و برای تحلیل و گزارش‌گیری مناسب‌تر است.

\section{لایه تجمیع}
لایه تجمیع، رخدادهای خام را به معیارهای قابل تحلیل تبدیل می‌کند. در این سامانه، تجمیع در پنجره‌های زمانی ده‌ثانیه‌ای انجام می‌شود. برای هر پنجره، رخدادها بر اساس کلیدهایی مانند شناسه فرایند، منبع رخداد، دسته رخداد و نوع رخداد گروه‌بندی می‌شوند و سپس معیارهایی مانند تعداد رخداد، مجموع مدت، میانگین مدت، بیشینه مدت و معیارهای اختصاصی هر دسته محاسبه می‌گردد.

وجود این لایه از چند جهت اهمیت دارد. نخست، تحلیل ناهنجاری معمولاً بر اساس روندها و معیارهای خلاصه‌شده معنادارتر از تحلیل تک‌تک رخدادهای خام است. دوم، تجمیع باعث کاهش حجم داده‌ای می‌شود که باید برای مدت طولانی‌تر نگه‌داری شود. در نتیجه، لازم نیست برای لایه جمع‌آوری داده، زمان نگه‌داری یا \lr{retention} زیادی در نظر گرفته شود. داده خام می‌تواند با زمان نگه‌داری کوتاه‌تر ذخیره شود و داده تجمیع‌شده، که حجم کم‌تر و ارزش تحلیلی بیش‌تری دارد، برای تحلیل و گزارش‌گیری حفظ گردد. این تصمیم باعث کاهش مصرف دیسک و ساده‌تر شدن مدیریت داده می‌شود.

سوم، لایه تجمیع مرزی میان مشاهده رخداد و تحلیل رفتاری ایجاد می‌کند. لایه جمع‌آوری صرفاً رخداد را ثبت می‌کند، اما لایه تجمیع آن رخدادها را به شاخص‌هایی تبدیل می‌کند که برای تصمیم‌گیری قابل استفاده‌اند.

\section{لایه تحلیل}
لایه تحلیل، مصرف‌کننده داده‌های تجمیع‌شده است. این لایه خروجی پنجره‌های زمانی را دریافت می‌کند و بررسی می‌کند که آیا معیارهای مشاهده‌شده نشانه‌ای از رفتار غیرعادی دارند یا خیر. از آنجا که ورودی این لایه داده خام نیست، تحلیل‌گر با معیارهایی کار می‌کند که از قبل خلاصه، دسته‌بندی و زمان‌بندی شده‌اند.

در این معماری، تحلیل ناهنجاری می‌تواند برای دسته‌های مختلف رخداد به صورت جداگانه انجام شود؛ برای نمونه پردازنده، حافظه، ورودی/خروجی فایل، شبکه، قفل‌ها و رخدادهای هسته. این جداسازی باعث می‌شود منطق تشخیص برای هر دسته متناسب با ماهیت همان دسته تعریف شود. برای مثال، معیار مناسب برای فشار پردازنده با معیار مناسب برای ازدحام قفل یا فشار حافظه یکسان نیست.

خروجی لایه تحلیل، گزارش رخدادهای تجمیع‌شده و گزارش ناهنجاری‌ها است. گزارش ناهنجاری نشان می‌دهد در چه بازه زمانی، برای چه دسته‌ای از رخدادها و بر اساس کدام معیار، رفتار غیرعادی مشاهده شده است.

\section{لایه ارزیابی}
در کنار لایه‌های اصلی سامانه، یک لایه ارزیابی نیز وجود دارد. این لایه بخشی از مسیر عادی پایش نیست، بلکه برای سنجش توان سامانه استفاده می‌شود. نقش آن ایجاد سناریوهای کنترل‌شده و بررسی این موضوع است که آیا لایه تحلیل می‌تواند ناهنجاری مورد انتظار را تشخیص دهد یا خیر.

لایه ارزیابی سناریوهایی مانند افزایش مصرف پردازنده، فشار حافظه و جمع‌آوری زباله، ورودی/خروجی سنگین فایل، ترافیک شبکه، رقابت روی قفل‌ها و سناریوهای وابسته به رخدادهای هسته را ایجاد می‌کند. سپس خروجی تحلیل‌گر بررسی می‌شود و گزارشی شامل وضعیت تشخیص، معیار تشخیص‌داده‌شده و تأخیر تشخیص تولید می‌گردد.

\section{مدل داده مشترک}
یکی از نکات کلیدی معماری، تعریف مدل داده مشترک پس از لایه جمع‌آوری است. هر رخداد، صرف نظر از این که از \lr{JFR} آمده باشد یا از مسیر هسته، باید اطلاعات پایه‌ای مشابهی داشته باشد. این اطلاعات شامل منبع رخداد، نوع رخداد، دسته رخداد، زمان، شناسه فرایند و داده‌های تکمیلی است.

وجود مدل مشترک باعث می‌شود لایه تجمیع بتواند هر دو مسیر داده را با منطق یکسان پردازش کند. همچنین لایه تحلیل می‌تواند بر اساس دسته و معیار تصمیم بگیرد، نه بر اساس جزئیات منبع تولید رخداد. این انتزاع یکی از دلایل توسعه‌پذیری معماری است؛ زیرا با اضافه شدن منبع داده جدید، کافی است خروجی آن منبع به مدل مشترک تبدیل شود.

\section{دلایل انتخاب معماری لایه‌ای}
انتخاب معماری لایه‌ای برای این سامانه چند دلیل اصلی دارد:
\begin{itemize}
\item \textbf{جداسازی مسئولیت‌ها:} هر لایه یک وظیفه مشخص دارد؛ جمع‌آوری، انتقال، تجمیع، تحلیل یا ارزیابی.
\item \textbf{کاهش وابستگی اجزا:} ارتباط اجزا از طریق جریان‌های \lr{Kafka} انجام می‌شود و هر بخش می‌تواند مستقل‌تر توسعه یا اجرا شود.
\item \textbf{مدیریت بهتر حجم داده:} داده خام با نگه‌داری کوتاه‌تر ذخیره می‌شود و داده تجمیع‌شده برای تحلیل و گزارش حفظ می‌گردد.
\item \textbf{پشتیبانی از چند منبع مشاهده:} معماری می‌تواند هم رخدادهای درون \lr{JVM} و هم رخدادهای سطح هسته را در یک مدل تحلیلی مشترک قرار دهد.
\item \textbf{امکان توسعه آینده:} افزودن منبع داده، معیار تجمیع یا قانون تحلیل جدید بدون تغییر بنیادین در کل سامانه امکان‌پذیر است.
\end{itemize}

\section{جمع‌بندی}
در این فصل معماری سامانه از دید لایه‌ای بررسی شد. سامانه دو مسیر داده اصلی دارد: مسیر رخدادهای \lr{JFR} و مسیر رخدادهای هسته. هر دو مسیر در لایه جمع‌آوری نرمال‌سازی شده و در \lr{Kafka} ذخیره می‌شوند. سپس لایه تجمیع، داده‌های خام را در پنجره‌های ده‌ثانیه‌ای به معیارهای خلاصه تبدیل می‌کند. لایه تحلیل این معیارها را برای تشخیص ناهنجاری مصرف می‌کند و لایه ارزیابی نیز با ایجاد سناریوهای کنترل‌شده، توان سامانه را می‌سنجد. این معماری با جداسازی مسئولیت‌ها و استفاده از مدل داده مشترک، امکان تحلیل بی‌درنگ و توسعه‌پذیر رخدادهای کارایی فرایندهای \lr{JVM} را فراهم می‌کند.
